\chapter{Dessins d'Enfants: The Genus of a Generation}
\label{ch:dessins}
% ═══════════════════════════════════════════════════════════════════════════════
% CENTRAL CORRESPONDENCE:
%   Discrete: Bipartite graph K_{2,N} with rotation system
%   Continuous: Riemann surface of genus g (Belyi)
%   The causal phase selects the rotation system → fixes the genus

\begin{introduction}
  \item The bridge problem: from the discrete graph $K_{2,N}$ to the
        continuous Riemann surface.
  \item Grothendieck's Dessins d'Enfants: every bipartite graph
        \textbf{is} a map on a surface.
  \item The monodromy twist: how Kuratowski phase entanglement
        forces a non-trivial genus.
  \item Generation--Genus correspondence: Gen~1 → sphere ($g=0$),
        Gen~2 → torus ($g=1$), Gen~3 → double torus ($g=2$).
\end{introduction}

% ─────────────────────────────────────────────────────────────────────────────
\section{The Bridge Problem}
% ─────────────────────────────────────────────────────────────────────────────

The preceding chapters established that matter consists of
Causal Prisms---bipartite subgraphs $K_{2,N}$ of the Hasse diagram.
Yang--Mills theory, the mathematical language of the Standard Model,
lives however on \emph{continuous manifolds}: bundles over a smooth
spacetime. How do we cross the bridge from a finite combinatorial graph
to a continuous Riemann surface?

\begin{figure}[H]
\centering
\begin{tikzpicture}[
  pole/.style={circle, draw=structure, fill=structure, minimum size=7pt,
               inner sep=0pt},
  belly/.style={circle, draw=main, fill=main!20, minimum size=7pt,
                inner sep=0pt},
  edge/.style={thick, structure!70},
  >=Stealth
]
% --- Poles ---
\node[pole, label={[color=structure]left:$u$ (past)}] (u) at (0,0) {};
\node[pole, label={[color=structure]left:$v$ (future)}] (v) at (0,4.5) {};

% --- Belly nodes ---
\node[belly, label={[color=main]right:$w_1$}] (w1) at (-1.8,2.25) {};
\node[belly, label={[color=main]right:$w_2$}] (w2) at (-0.6,2.25) {};
\node[belly, label={[color=main]left:$\cdots$}]  (wd) at ( 0.6,2.25) {};
\node[belly, label={[color=main]left:$w_n$}]  (wn) at ( 1.8,2.25) {};

% --- Edges from u ---
\draw[edge, decoration={markings, mark=at position 0.55 with
      {\arrow{>}}}, postaction={decorate}] (u) -- (w1)
      node[midway, left, font=\scriptsize, color=structure!60] {$e_1$};
\draw[edge, decoration={markings, mark=at position 0.55 with
      {\arrow{>}}}, postaction={decorate}] (u) -- (w2)
      node[midway, left, font=\scriptsize, color=structure!60] {$e_2$};
\draw[edge, densely dotted] (u) -- (wd);
\draw[edge, decoration={markings, mark=at position 0.55 with
      {\arrow{>}}}, postaction={decorate}] (u) -- (wn)
      node[midway, right, font=\scriptsize, color=structure!60] {$e_n$};

% --- Edges from belly to v ---
\draw[edge, decoration={markings, mark=at position 0.55 with
      {\arrow{>}}}, postaction={decorate}] (w1) -- (v)
      node[pos=0.4, left, font=\scriptsize, color=structure!60] {$e_{n+1}$};
\draw[edge, decoration={markings, mark=at position 0.55 with
      {\arrow{>}}}, postaction={decorate}] (w2) -- (v)
      node[pos=0.3, right, font=\scriptsize, color=structure!60] {$e_{n+2}$};
\draw[edge, densely dotted] (wd) -- (v);
\draw[edge, decoration={markings, mark=at position 0.55 with
      {\arrow{>}}}, postaction={decorate}] (wn) -- (v)
      node[pos=0.4, right, font=\scriptsize, color=structure!60] {$e_{2n}$};

% --- Time arrow ---
\draw[->, thick, structure!40] (2.8,0) -- (2.8,4.5)
      node[midway, right, font=\small, color=structure!50] {time};

% --- Brace for belly ---
\draw[decorate, decoration={brace, amplitude=5pt, mirror},
      main!70] (-2.1,1.7) -- (2.1,1.7)
      node[midway, below=6pt, font=\small, color=main] {belly: $N$ intermediaries};
\end{tikzpicture}
\caption{The Causal Prism $K_{2,N}$. The black vertices $u$ (past
pole) and $v$ (future pole) define the time direction. The golden
vertices $w_i$ form the \emph{belly}---the internal structure of the
prism. Each directed edge is a causal link of the Hasse diagram.
The graph is bipartite: only pole$\to$intermediary and
intermediary$\to$pole edges exist, never pole$\to$pole nor
intermediary$\to$intermediary.}
\label{fig:prisma-base}
\end{figure}

The answer was discovered by Grothendieck~\cite{grothendieck1984}:
every bipartite graph, equipped with a \emph{cyclic order} of edges
at each vertex, \textbf{is} a cellular map on a compact oriented
surface. The \textbf{genus} of that surface---the number of
``handles'' or ``holes''---depends not on the abstract graph (which for
$K_{2,N}$ is always planar), but on how the edges are cyclically
ordered at the vertices. This correspondence is called a
\textbf{dessin d'enfant} (``child's drawing'').

% ─────────────────────────────────────────────────────────────────────────────
\section{The Dessin Formalism}
\label{sec:dessins-formalismo}
% ─────────────────────────────────────────────────────────────────────────────

\begin{definition}[Dessin d'Enfant]
\label{def:dessin}
A \textbf{dessin d'enfant} is a triple $(\sigma_0, \sigma_1, E)$
where:
\begin{itemize}
  \item $E = \{1, 2, \ldots, m\}$ is the \textbf{edge set},
  \item $\sigma_0 \in S_m$ is a permutation whose cycles encode the
        \textbf{cyclic order of edges around each black vertex},
  \item $\sigma_1 \in S_m$ is a permutation whose cycles encode the
        \textbf{cyclic order of edges around each white vertex}.
\end{itemize}
The group $\langle \sigma_0, \sigma_1 \rangle \le S_m$ acts
transitively on $E$ if and only if the underlying graph is connected.
\end{definition}

\begin{definition}[Face Permutation]
\label{def:perm-caras}
The \textbf{face permutation} is:
\[
  \sigma_\infty \;=\; (\sigma_0 \circ \sigma_1)^{-1},
\]
where $(\sigma_0 \circ \sigma_1)(e) = \sigma_0(\sigma_1(e))$. Each
cycle of $\sigma_\infty$ traces the boundary of exactly one
\textbf{face} of the dessin. The number of faces is:
\[
  F \;=\; \#\{\text{cycles of } \sigma_\infty\}
    \;=\; \#\{\text{cycles of } \sigma_0 \circ \sigma_1\},
\]
since inversion preserves the cycle structure.
\end{definition}

\begin{theorem}[Euler--Grothendieck Formula]
\label{thm:euler-grothendieck}
Let $(\sigma_0, \sigma_1, E)$ be a connected dessin. We define:
\begin{align*}
  V &= \#\{\text{cycles of } \sigma_0\}
     + \#\{\text{cycles of } \sigma_1\}
     && \text{(vertices)}, \\
  E_{\mathrm{total}} &= |E|
     && \text{(edges)}, \\
  F &= \#\{\text{cycles of } \sigma_\infty\}
     && \text{(faces)}.
\end{align*}
Then $V - E_{\mathrm{total}} + F = 2 - 2g$, where $g \ge 0$ is the
\textbf{genus} of the unique compact oriented surface on which
the dessin embeds cellularly.
\end{theorem}

\begin{proof}
This is the Euler--Poincar\'e formula for the cellular decomposition of a
compact oriented surface: 0-cells (vertices), 1-cells (edges)
and 2-cells (faces bounded by the cycles of $\sigma_\infty$). The
Euler characteristic $\chi = V - E + F$ determines the genus via
$\chi = 2 - 2g$. The existence and uniqueness of such a decomposition for
a given rotation system is constructive: see
Lando--Zvonkin~\cite{lando2004}, Chapter~1.
\end{proof}

\begin{theorem}[Belyi's Theorem]
\label{thm:belyi}
Every compact Riemann surface $X$ defined over
$\overline{\mathbb{Q}}$ admits a Belyi function
$\beta \colon X \to \mathbb{P}^1$ ramified only over
$\{0, 1, \infty\}$, whose preimage of $[0,1]$ is a dessin.
Conversely, every dessin determines a unique Riemann surface and
Belyi function up to isomorphism.
\end{theorem}

\noindent
Belyi's theorem is the mathematical miracle that makes the bridge possible:
a purely combinatorial object (a pair of permutations) determines a
unique algebraic curve. The physical content of this chapter is that the
\emph{causal structure} of the prism selects the pair of permutations.

% ─────────────────────────────────────────────────────────────────────────────
\section{Generation 1: The Sphere ($g = 0$)}
\label{sec:gen1-esfera}
% ─────────────────────────────────────────────────────────────────────────────

Before tackling the physically crucial case, let us establish the method
with the simplest example.

\subsection{$K_{2,3}$ without Entanglement}

Consider $K_{2,3}$ with black vertices $u, v$ (poles) and white
vertices $w_1, w_2, w_3$ (intermediaries). We label the 6 edges:
\[
  \underbrace{e_1 = (u,w_1),\; e_2 = (u,w_2),\; e_3 = (u,w_3)}_{\text{from } u}
  \qquad
  \underbrace{e_4 = (v,w_1),\; e_5 = (v,w_2),\; e_6 = (v,w_3)}_{\text{from } v}
\]
Each white vertex has degree~2, so $\sigma_1$ is forced:
\[
  \sigma_1 = (1\;4)(2\;5)(3\;6).
\]

\noindent\textbf{The causal antipodal reversal.} In a Causal Prism
embedded in Minkowski $3{+}1$, the celestial spheres of the past
pole $u$ and future pole $v$ are related by the antipodal map of $S^2$,
which \textbf{reverses the orientation}. If from $u$ the intermediaries are
seen in cyclic order $(w_1, w_2, w_3)$, from $v$ they are seen in the
reverse order $(w_1, w_3, w_2)$. In terms of edges:
\[
  \sigma_0 = (1\;2\;3)(4\;6\;5).
\]

\noindent
Let us compute $\sigma_0(\sigma_1(e))$ for each edge:
\begin{center}
\renewcommand{\arraystretch}{1.15}
\begin{tabular}{c@{\;\;}c@{\;\;}c@{\;\;}c@{\;\;}c@{\;\;}c@{\;\;}c}
\toprule
$e$ & 1 & 2 & 3 & 4 & 5 & 6 \\
\midrule
$\sigma_1(e)$ & 4 & 5 & 6 & 1 & 2 & 3 \\
$\sigma_0(\sigma_1(e))$ & 6 & 4 & 5 & 2 & 3 & 1 \\
\bottomrule
\end{tabular}
\end{center}
The cycles of $\sigma_0 \circ \sigma_1$: $(1\;6)(2\;4)(3\;5)$---three
2-cycles, so $F = 3$ faces.

\noindent
\textbf{Verification:}\quad $V = 2 + 3 = 5$,\quad $E = 6$,\quad $F = 3$.

\[
  V - E + F = 5 - 6 + 3 = 2 = 2 - 2(0)
  \quad\Longrightarrow\quad \boxed{g = 0} \text{ (sphere).}
\]

\begin{figure}[H]
\centering
\begin{tikzpicture}[
  pole/.style={circle, draw=structure, fill=structure, minimum size=7pt,
               inner sep=0pt},
  belly/.style={circle, draw=main, fill=main!20, minimum size=7pt,
                inner sep=0pt},
  edge/.style={thick},
  face/.style={fill opacity=0.12, draw=none},
  >=Stealth
]
% --- Poles ---
\node[pole, label={[color=structure]below:$u$}] (u) at (0,0) {};
\node[pole, label={[color=structure]above:$v$}] (v) at (0,4) {};

% --- Belly nodes ---
\node[belly, label={[color=main]left:$w_1$}]  (w1) at (-2,2) {};
\node[belly, label={[color=main]right:$w_2$}] (w2) at ( 0,2) {};
\node[belly, label={[color=main]right:$w_3$}] (w3) at ( 2,2) {};

% --- Faces (shaded regions) ---
\fill[fourth, face] (u.center) -- (w1.center) -- (v.center) -- cycle;
\fill[second, face] (u.center) -- (w2.center) -- (v.center) -- cycle;
\fill[third, face]  (u.center) -- (w3.center) -- (v.center) -- cycle;

% --- Edges ---
\draw[edge, fourth!80!black] (u) -- (w1)
  node[midway, left, font=\scriptsize] {$e_1$};
\draw[edge, second!80!black] (u) -- (w2)
  node[midway, left, font=\scriptsize] {$e_2$};
\draw[edge, third!80!black]  (u) -- (w3)
  node[midway, right, font=\scriptsize] {$e_3$};
\draw[edge, fourth!80!black] (w1) -- (v)
  node[midway, left, font=\scriptsize] {$e_4$};
\draw[edge, second!80!black] (w2) -- (v)
  node[midway, left, font=\scriptsize] {$e_5$};
\draw[edge, third!80!black]  (w3) -- (v)
  node[midway, right, font=\scriptsize] {$e_6$};

% --- Face labels ---
\node[font=\scriptsize, fourth!80!black] at (-1.15,2) {$f_1$};
\node[font=\scriptsize, second!80!black] at (-0.15,2.6) {$f_2$};
\node[font=\scriptsize, third!80!black]  at ( 1.15,2) {$f_3$};

% --- Sphere annotation ---
\draw[<-, structure!50, thick] (3,2) -- (4.2,2)
  node[right, font=\small, text width=3.2cm, color=structure!70]
  {$F = 3$ faces\\$g = 0$ (sphere)\\No crossings};
\end{tikzpicture}
\caption{Generation~1: the dessin of $K_{2,3}$ on the sphere.
Each pair of conjugate edges $(e_i, e_{i+3})$ encloses an
independent face~$f_i$. The three causal channels do not cross: the graph
is planar and embeds seamlessly on $S^2$. This is the electron.}
\label{fig:gen1-esfera}
\end{figure}

\noindent
Each face surrounds exactly one pair of conjugate edges
$(e_i, e_{i+3})$ connecting the two poles to the same intermediary.
There is no entanglement: each causal path closes cleanly
upon itself. The sphere has no holes.

\subsection{The General Formula for Non-Entangled Prisms}

For any $K_{2,N}$ without phase entanglement:
$V = N + 2$, $E = 2N$, $F = N$ (since each intermediary generates an
independent 2-cycle face). Then:
\[
  (N+2) - 2N + N = 2 = 2 - 2g
  \quad\Longrightarrow\quad g = 0 \text{ for all } N.
\]
A non-entangled prism is \textbf{always a sphere},
regardless of its topological mass.

% ─────────────────────────────────────────────────────────────────────────────
\section{Generation 2: The Torus ($g = 1$)}
\label{sec:gen2-toro}
% ─────────────────────────────────────────────────────────────────────────────

This is the central result of the chapter: the rigorous proof that
Kuratowski phase entanglement forces the Riemann surface to tear itself
apart and reform into a torus.

\subsection{Setup: 8 Edges of $K_{2,4}$}

We have 2 poles (black: past $u$, future $v$) and 4 intermediaries
(white: $w_1, w_2, w_3, w_4$). We label the 8 edges:
\[
  \underbrace{e_1 = (u,w_1),\; e_2 = (u,w_2),\; e_3 = (u,w_3),\;
    e_4 = (u,w_4)}_{\text{from } u}
  \quad
  \underbrace{e_5 = (v,w_1),\; e_6 = (v,w_2),\; e_7 = (v,w_3),\;
    e_8 = (v,w_4)}_{\text{from } v}
\]

\subsection{The Intermediary Permutation ($\sigma_1$)}

Since each intermediary $w_i$ connects only to the two poles, its
permutation is the deterministic swap of its two edges:
\[
  \sigma_1 = (1\;5)(2\;6)(3\;7)(4\;8).
\]

\subsection{The Polar Permutation ($\sigma_0$) with the Kuratowski Twist}
\label{subsec:giro-kuratowski}

This is where the physics intervenes.

If this were an isolated graph, the edges would form
a simple planar ring. From $u$'s perspective, the cyclic order
is $(1\;2\;3\;4)$. Due to the antipodal reversal, $v$ would see
exactly the reverse: $(8\;7\;6\;5)$. If we compute the faces
for that configuration, we obtain $F = 4$ and $g = 0$.

But a $K_{2,4}$ \textbf{does not exist in isolation}: it is embedded in the
dense Hasse diagram of the causal set. The following observation
is crucial.

\begin{lemma}[Ambient-Induced Non-Planarity]
\label{lem:ambient-nonplanarity}
Let $\Pi = K_{2,4}$ be a Causal Prism embedded in a Hasse diagram
$\mathcal{H}$ of density $\rho \gg 1$. Although $K_{2,4}$ is
intrinsically planar (every complete bipartite graph $K_{2,N}$ admits
a planar embedding for any~$n$), the ``straight'' planar embedding
(without crossings) within~$\mathcal{H}$ is
\textbf{forbidden} by Theorem~\ref{thm:kuratowski-threat} (Kuratowski
Threat).

Specifically: in the straight embedding, the vacuum nodes surrounding
the prism share links with both poles $u, v$ and with at least
$K_{\text{threat}} = 2$ intermediaries. This forms $K_5$ minors
between the vacuum nodes and the prism. To avoid the planarity
violation, the rotation system of the prism must \textbf{deform
its cyclic order} at least at one pole, swapping the position
of two adjacent edges. This deformation is the \textbf{Kuratowski
twist}.
\end{lemma}

\begin{proof}
Let $t$ be a vacuum node exterior to the prism with $t \sim u$,
$t \sim v$, and $t \sim w_i$ for at least $K_{\text{threat}} = 2$
intermediaries. If the prism's edges preserve their planar cyclic order,
the subgraph induced by $\{u, v, w_i, w_j, t\}$ contains a
$K_5$ minor: all 5~nodes are mutually connected (directly or
transitively through Hasse links). Theorem~\ref{thm:kuratowski-threat}
forbids this configuration. The
only resolution that preserves the global planarity
of~$\mathcal{H}$ without absorbing~$t$ into a pole (which would destroy the
prism) is to modify the cyclic order of $\sigma_0(v)$,
introducing a crossing between the edges of the intermediaries that
interact with~$t$. In a Poisson diamond with $\rho \gg 1$,
the probability that at least one such node~$t$ exists tends to~1
for prisms with $N \geq 4$.
\end{proof}

Two of the internal causal paths---say $w_2$ and $w_3$---are
forced to cross in the topological embedding, not by an
intrinsic property of $K_{2,4}$, but by the planarity pressure
exerted by the ambient Hasse diagram
(Lemma~\ref{lem:ambient-nonplanarity}).

Because of this crossing, the future pole $v$ does \textbf{not} see a clean
inverted ring. It sees a twisted cyclic order where edges
$e_6$ and $e_7$ have swapped positions:

\medskip
\noindent
\textbf{At $u$:}\quad $\sigma_0(u) = (1\;2\;3\;4)$

\noindent
\textbf{At $v$:}\quad $\sigma_0(v) = (8\;\mathbf{6\;7}\;5)$
\quad\textit{(6 and 7 swapped relative to the clean reversal
$(8\;7\;6\;5)$)}

\medskip
\noindent
The combined polar permutation is:
\[
  \sigma_0 = (1\;2\;3\;4)(8\;6\;7\;5).
\]

\subsection{The Face Computation ($\sigma_\infty$)}

In Grothendieck's theory of Dessins d'Enfants~\cite{grothendieck1984},
the continuous faces of the Riemann surface are strictly defined by the
cycles of $\sigma_0 \circ \sigma_1$. Let us trace the composition
$\sigma_0(\sigma_1(e))$ step by step.

\medskip
\noindent
\textbf{Cycle 1 --- The non-entangled path:}
\[
  1 \xrightarrow{\;\sigma_1\;} 5 \xrightarrow{\;\sigma_0\;} 8
  \qquad\qquad
  8 \xrightarrow{\;\sigma_1\;} 4 \xrightarrow{\;\sigma_0\;} 1
\]
\textbf{First face:} $(1\;8)$---a clean 2-cycle. Edges $e_1$
and $e_8$ (connecting the poles to $w_1$ and $w_4$ respectively)
close into an independent face. No entanglement.

\medskip
\noindent
\textbf{Cycle 2 --- The entangled compression:}
\begin{align*}
  2 &\xrightarrow{\;\sigma_1\;} 6 \xrightarrow{\;\sigma_0\;} 7 \\
  7 &\xrightarrow{\;\sigma_1\;} 3 \xrightarrow{\;\sigma_0\;} 4 \\
  4 &\xrightarrow{\;\sigma_1\;} 8 \xrightarrow{\;\sigma_0\;} 6 \\
  6 &\xrightarrow{\;\sigma_1\;} 2 \xrightarrow{\;\sigma_0\;} 3 \\
  3 &\xrightarrow{\;\sigma_1\;} 7 \xrightarrow{\;\sigma_0\;} 5 \\
  5 &\xrightarrow{\;\sigma_1\;} 1 \xrightarrow{\;\sigma_0\;} 2
\end{align*}
\textbf{Second face:} $(2\;7\;4\;6\;3\;5)$---a single massive
6-cycle. The crossing of just two discrete paths has collapsed what
should have been three independent cycles into a single entangled
orbit of 6 edges.

\medskip
\noindent
Summary table:
\begin{center}
\renewcommand{\arraystretch}{1.15}
\begin{tabular}{c@{\;\;}c@{\;\;}c@{\;\;}c@{\;\;}c@{\;\;}c@{\;\;}c@{\;\;}c@{\;\;}c}
\toprule
$e$ & 1 & 2 & 3 & 4 & 5 & 6 & 7 & 8 \\
\midrule
$\sigma_1(e)$ & 5 & 6 & 7 & 8 & 1 & 2 & 3 & 4 \\
$\sigma_0(\sigma_1(e))$ & 8 & 7 & 4 & 6 & 2 & 3 & 5 & 1 \\
\bottomrule
\end{tabular}
\end{center}
$\sigma_0 \circ \sigma_1 = (1\;8)(2\;7\;4\;6\;3\;5)$
\quad$\Longrightarrow$\quad $F = 2$ faces.

\subsection{The Topological Verdict}

Let us observe what the Kuratowski pressure did to the mathematics.
The crossing of just two discrete paths collapsed what should have
been three independent cycles into a single orbit of 6 edges.

\[
  V - E + F = 6 - 8 + 2 = 0 = 2 - 2g
  \quad\Longrightarrow\quad \boxed{g = 1} \text{ (torus).}
\]
Each twist reduces~$F$ by exactly~2 at fixed $V$ and~$E$: from
$\chi = V - E + F = 2 - 2g$ it follows that $\Delta F = -2\,\Delta g$.

\begin{figure}[H]
\centering
\begin{tikzpicture}[
  pole/.style={circle, draw=structure, fill=structure, minimum size=7pt,
               inner sep=0pt},
  belly/.style={circle, draw=main, fill=main!20, minimum size=7pt,
                inner sep=0pt},
  edge/.style={thick},
  >=Stealth
]
% ======= LEFT PANEL: planar K_{2,4} with crossing =======
\begin{scope}[shift={(-4,0)}]
  \node[font=\small\bfseries, color=structure] at (0,5) {The discrete crossing};

  % Poles
  \node[pole, label={[color=structure]below:$u$}] (u1) at (0,0) {};
  \node[pole, label={[color=structure]above:$v$}] (v1) at (0,4) {};

  % Belly
  \node[belly, label={[color=main]left:$w_1$}]  (a1) at (-2.1,2) {};
  \node[belly, label={[color=main]left:$w_2$}]  (a2) at (-0.7,2) {};
  \node[belly, label={[color=main]right:$w_3$}] (a3) at ( 0.7,2) {};
  \node[belly, label={[color=main]right:$w_4$}] (a4) at ( 2.1,2) {};

  % Clean edges (outer, no crossing)
  \draw[edge, fourth!80!black] (u1) -- (a1)
    node[midway, left, font=\scriptsize] {$e_1$};
  \draw[edge, fourth!80!black] (a1) -- (v1)
    node[midway, left, font=\scriptsize] {$e_5$};
  \draw[edge, fourth!80!black] (u1) -- (a4)
    node[midway, right, font=\scriptsize] {$e_4$};
  \draw[edge, fourth!80!black] (a4) -- (v1)
    node[midway, right, font=\scriptsize] {$e_8$};

  % From u: straight edges to w2, w3 (no crossing below belly)
  \draw[edge, third!80!black] (u1) -- (a2)
    node[midway, left, font=\scriptsize] {$e_2$};
  \draw[edge, second!80!black] (u1) -- (a3)
    node[midway, right, font=\scriptsize] {$e_3$};

  % From v: CROSSED — e6 (v→w2) curves RIGHT, e7 (v→w3) curves LEFT
  % so they physically cross each other
  % e7: v→w3, exits v heading LEFT, arrives at w3 on the RIGHT
  \draw[edge, second!80!black] (v1) to[out=-110, in=90] (a3);
  % e6: v→w2, exits v heading RIGHT, arrives at w2 on the LEFT
  \draw[edge, third!80!black] (v1) to[out=-70, in=90] (a2);

  % Labels placed manually near v
  \node[font=\scriptsize, second!80!black] at (-0.55,3.55) {$e_7$};
  \node[font=\scriptsize, third!80!black]  at ( 0.55,3.55) {$e_6$};

  % Crossing highlight
  \node[circle, draw=third, very thick, inner sep=3pt,
        fill=white] at (0,2.95) {};
  \node[font=\bfseries\small, color=third!80!black] at (0,2.95) {$\times$};

  % Annotation
  \draw[->, thick, third!60!black] (0.35,3.15) -- (1.1,3.6)
    node[right, font=\scriptsize, text width=2.2cm, color=third!70!black]
    {Kuratowski\\[-2pt]twist:\\$e_6 \leftrightarrow e_7$};
\end{scope}

% ======= RIGHT PANEL: the torus =======
\begin{scope}[shift={(4.5,2)}]
  \node[font=\small\bfseries, color=structure] at (0,3) {The continuous surface};

  % Torus as fundamental polygon
  \draw[very thick, structure] (-2,-1.8) rectangle (2,1.8);

  % Identification arrows on edges
  \draw[very thick, third, decoration={markings,
    mark=at position 0.5 with {\arrow{>}}},
    postaction={decorate}] (-2,-1.8) -- (2,-1.8);
  \draw[very thick, third, decoration={markings,
    mark=at position 0.5 with {\arrow{>}}},
    postaction={decorate}] (-2,1.8) -- (2,1.8);
  \draw[very thick, second, decoration={markings,
    mark=at position 0.5 with {\arrow{>>}}},
    postaction={decorate}] (-2,-1.8) -- (-2,1.8);
  \draw[very thick, second, decoration={markings,
    mark=at position 0.5 with {\arrow{>>}}},
    postaction={decorate}] (2,-1.8) -- (2,1.8);

  % Interior: the two faces
  \fill[fourth, opacity=0.10] (-2,-1.8) rectangle (2,1.8);

  % Poles on the torus
  \node[pole] (p1) at (-1.3,0) {};
  \node[font=\scriptsize, color=structure, below=2pt] at (p1) {$u$};
  \node[pole] (p2) at (1.3,0) {};
  \node[font=\scriptsize, color=structure, below=2pt] at (p2) {$v$};

  % Face labels
  \node[font=\small, main] at (0,0.9) {small face};
  \node[font=\small, main] at (0,-0.9) {large face (6-cycle)};

  % Annotation
  \node[font=\scriptsize, color=structure!70, text width=3.4cm, align=center]
    at (0,-2.6) {Fundamental polygon of the torus:\\identify opposite sides.\\
    $F = 2$,\; $g = 1$.};
\end{scope}
\end{tikzpicture}
\caption{Generation~2: the Kuratowski twist and the torus.
\textbf{Left:} the graph $K_{2,4}$ with the forced crossing between
$e_6$ and $e_7$ --- the phase anti-correlation forces $w_2$ and $w_3$ to
repel each other. This crossing cannot be resolved in the plane.
\textbf{Right:} the fundamental polygon of the torus ($T^2$) that resolves
the crossing: by identifying opposite sides, the edges separate without
cutting each other. The surface has only $F = 2$ faces and genus $g = 1$.
This is the muon.}
\label{fig:gen2-toro}
\end{figure}

\begin{result}[title={The rigorous proof of topological mass}]
One does not ``assert'' that $F = 2$. It is \textbf{proved} via explicit
permutation arithmetic: Kuratowski anti-correlation $\to$
monodromy twist $\to$ cycle fusion $\to$ higher genus.

The topological inertia (mass) of the muon is greater than the electron's
because a random walker entering through path~2 does not simply
pass toward the future: it is forced by the permutation $\sigma_\infty$ to
traverse an entangled orbit of 6 steps
$(2 \to 7 \to 4 \to 6 \to 3 \to 5)$ before it can resolve. The
geometry of the continuum has literally folded into a torus to
accommodate the crossing.
\end{result}

\stoneref{The monodromy twist is the \textbf{irreducibility of a
shared lemma}. On the sphere ($g = 0$), every deductive path contracts
to a point: there is a single logical route from past to future. On
the torus ($g = 1$), the crossing of two phase paths creates a
\textbf{non-contractible} deductive loop---a proof that cannot be
simplified to any other. The genus counts the irreducible
proofs in the Heyting algebra of the prism.}

% ─────────────────────────────────────────────────────────────────────────────
\section{Generation 3: The Double Torus ($g = 2$)}
\label{sec:gen3-doble-toro}
% ─────────────────────────────────────────────────────────────────────────────

A Generation~3 particle possesses three phase classes
$\{+, 0, -\}$ in its belly. Each pair of adjacent classes generates an
independent Kuratowski crossing. With two crossings,
Definition~\ref{def:perm-caras} yields:
\[
  F = N - 4 \qquad\Longrightarrow\qquad
  (N+2) - 2N + (N-4) = -2 = 2 - 2(2)
  \qquad\Longrightarrow\qquad g = 2.
\]

To verify, let us apply the mechanism to $K_{2,6}$ with three phase
classes $W_+ = \{w_1, w_2\}$, $W_0 = \{w_3, w_4\}$,
$W_- = \{w_5, w_6\}$. We label 12 edges: $e_1, \ldots, e_6$
from $u$ and $e_7, \ldots, e_{12}$ from $v$
($e_i$ and $e_{i+6}$ share intermediary $w_i$).

\[
  \sigma_1 = (1\;7)(2\;8)(3\;9)(4\;10)(5\;11)(6\;12).
\]

Without entanglement, the cyclic order at $u$ is $(1\;2\;3\;4\;5\;6)$ and
the antipodal reversal at $v$ is $(12\;11\;10\;9\;8\;7)$. This would give
$F = 6$ and $g = 0$. But Generation~3 requires \textbf{two} Kuratowski
twists to avoid the forbidden $K_5$ subgraphs:
\begin{itemize}
  \item \textbf{Twist 1:} The paths $w_2$ and $w_3$ repel/cross
        (edges $e_8$ and $e_9$ swap at $v$).
  \item \textbf{Twist 2:} The paths $w_4$ and $w_5$ repel/cross
        (edges $e_{10}$ and $e_{11}$ swap at $v$).
\end{itemize}
The future pole $v$ now sees a \textbf{doubly twisted} cyclic order:

\medskip
\noindent
\textbf{At $u$:}\quad $\sigma_0(u) = (1\;2\;3\;4\;5\;6)$

\noindent
\textbf{At $v$:}\quad $\sigma_0(v) = (12\;10\;11\;8\;9\;7)$
\quad\textit{(10$\leftrightarrow$11 and 8$\leftrightarrow$9 swapped
relative to the clean reversal $(12\;11\;10\;9\;8\;7)$)}

\medskip
\noindent
The combined polar permutation is:
\[
  \sigma_0 = (1\;2\;3\;4\;5\;6)(12\;10\;11\;8\;9\;7).
\]

\medskip
\noindent
We trace $\sigma_0(\sigma_1(e))$ step by step through the double
entanglement.

\medskip
\noindent
\textbf{Cycle 1 --- The outer boundary:}
\[
  1 \xrightarrow{\;\sigma_1\;} 7 \xrightarrow{\;\sigma_0\;} 12
  \qquad\qquad
  12 \xrightarrow{\;\sigma_1\;} 6 \xrightarrow{\;\sigma_0\;} 1
\]
\textbf{First face:} $(1\;12)$---a clean 2-cycle.

\medskip
\noindent
\textbf{Cycle 2 --- The double-twist orbit:}

Observe how the double crossing traps the remaining 10 edges in a
\textbf{single massive recursive loop}:
\begin{align*}
  2 &\xrightarrow{\;\sigma_1\;} 8  \xrightarrow{\;\sigma_0\;} 9 \\
  9 &\xrightarrow{\;\sigma_1\;} 3  \xrightarrow{\;\sigma_0\;} 4 \\
  4 &\xrightarrow{\;\sigma_1\;} 10 \xrightarrow{\;\sigma_0\;} 11 \\
  11 &\xrightarrow{\;\sigma_1\;} 5  \xrightarrow{\;\sigma_0\;} 6 \\
  6 &\xrightarrow{\;\sigma_1\;} 12 \xrightarrow{\;\sigma_0\;} 10
    \quad\textit{(twist~2 deflects it!)} \\
  10 &\xrightarrow{\;\sigma_1\;} 4  \xrightarrow{\;\sigma_0\;} 5 \\
  5 &\xrightarrow{\;\sigma_1\;} 11 \xrightarrow{\;\sigma_0\;} 8
    \quad\textit{(twist~1 deflects it!)} \\
  8 &\xrightarrow{\;\sigma_1\;} 2  \xrightarrow{\;\sigma_0\;} 3 \\
  3 &\xrightarrow{\;\sigma_1\;} 9  \xrightarrow{\;\sigma_0\;} 7 \\
  7 &\xrightarrow{\;\sigma_1\;} 1  \xrightarrow{\;\sigma_0\;} 2
    \quad\textit{(the orbit closes)}
\end{align*}
\textbf{Second face:}
$(2\;9\;4\;11\;6\;10\;5\;8\;3\;7)$---a single 10-cycle.

\medskip
\noindent
Summary table:
\begin{center}
\renewcommand{\arraystretch}{1.15}
\resizebox{\linewidth}{!}{%
\begin{tabular}{c@{\;\;}c@{\;\;}c@{\;\;}c@{\;\;}c@{\;\;}c@{\;\;}c@{\;\;}c@{\;\;}c@{\;\;}c@{\;\;}c@{\;\;}c@{\;\;}c}
\toprule
$e$ & 1 & 2 & 3 & 4 & 5 & 6 & 7 & 8 & 9 & 10 & 11 & 12 \\
\midrule
$\sigma_1(e)$
  & 7 & 8 & 9 & 10 & 11 & 12 & 1 & 2 & 3 & 4 & 5 & 6 \\
$\sigma_0(\sigma_1(e))$
  & 12 & 9 & 7 & 11 & 8 & 10 & 2 & 3 & 4 & 5 & 6 & 1 \\
\bottomrule
\end{tabular}}
\end{center}
$\sigma_0 \circ \sigma_1 = (1\;12)(2\;9\;4\;11\;6\;10\;5\;8\;3\;7)$
\quad$\Longrightarrow$\quad $F = 2$ faces.

\medskip
\noindent
\textbf{Verification:}
\[
  V - E + F = 8 - 12 + 2 = -2 = 2 - 2(2)
  \quad\Longrightarrow\quad \boxed{g = 2} \text{ (double torus).}
\]

\begin{result}[title={The physics of the Top quark}]
When a unit of energy interacts with a Generation~1 particle
($g = 0$), it slides over a smooth topological sphere.
When it interacts with a Generation~3 particle ($g = 2$), the graph
forces it to traverse a \textbf{double torus}---a space with two
fundamental holes.

The random walker gets trapped in that massive recursive orbit of
10 steps $(2 \to 9 \to 4 \to 11 \to 6 \to 10 \to 5 \to 8 \to 3 \to 7)$,
bouncing between the topological twists before it can exit through the
future pole. That severe geometric delay \textbf{is} the particle's
mass.
\end{result}

\noindent
The pattern is clear: each Kuratowski crossing absorbs independent
cycles into a larger entangled orbit, forcing the Riemann surface
to add a topological handle.

% ─────────────────────────────────────────────────────────────────────────────
\section{The Generation--Genus Correspondence}
\label{sec:gen-genero}
% ─────────────────────────────────────────────────────────────────────────────

\begin{theorem}[Generation--Genus Correspondence]
\label{thm:gen-genero}
The topological genus of the Riemann surface underlying a Causal Prism
of generation $g_{\mathrm{gen}}$ is:
\[
  g \;=\; g_{\mathrm{gen}} - 1.
\]
\end{theorem}

\begin{proof}
In each case, $V = N + 2$, $E = 2N$, and $\sigma_1$ consists of $N$
transpositions.

\medskip\noindent
\textbf{Case 1: Generation~1} ($g_{\mathrm{gen}} = 1$, a single phase
class).\\
No anti-correlation between classes. The rotation system is the
clean antipodal reversal. $F = N$ faces of 2-cycles:
\[
  (N+2) - 2N + N = 2 \;\Longrightarrow\; g = 0 \text{ (sphere).}
\]

\medskip\noindent
\textbf{Case 2: Generation~2} ($g_{\mathrm{gen}} = 2$, two phase
classes).\\
An entangled pair produces a Kuratowski crossing
(\S\ref{subsec:giro-kuratowski}). $F = N - 2$:
\[
  (N+2) - 2N + (N-2) = 0 \;\Longrightarrow\; g = 1 \text{ (torus).}
\]

\medskip\noindent
\textbf{Case 3: Generation~3} ($g_{\mathrm{gen}} = 3$, three phase
classes).\\
Two entangled pairs produce two independent crossings. $F = N - 4$:
\[
  (N+2) - 2N + (N-4) = -2 \;\Longrightarrow\; g = 2
  \text{ (double torus).}
\]
\end{proof}

\medskip
\begin{center}
\renewcommand{\arraystretch}{1.2}
\begin{tabular}{cclccc}
\toprule
Gen & $g$ & Surface & $\langle n \rangle$ & Particles
  & Crossings \\
\midrule
1 & 0 & Sphere        & 4.555 & $e,\; u,\; d$ & 0 \\
2 & 1 & Torus         & 6.531 & $\mu,\; c,\; s$ & 1 \\
3 & 2 & Double torus  & 7.732 & $\tau,\; t,\; b$ & 2 \\
\bottomrule
\end{tabular}
\end{center}

\medskip
\noindent
Since $g_{\mathrm{gen}} \le 3$ (the generation bound theorem),
$g \le 2$. No Causal Prism can live on a surface of genus~3 or higher.

The mass ratios $\langle n \rangle$ between generations---$1 :
1{.}434 : 1{.}698$---are not free parameters: they are derived
analytically from the belly distributions and the phase fractions
via a zero-parameter occupancy model. The complete derivation
is presented in \textbf{Practical Demonstration~III}
(p.~\pageref{sec:demo-mass}).

% ─────────────────────────────────────────────────────────────────────────────
\section{Module-Topological Compression}
\label{sec:compresion}
% ─────────────────────────────────────────────────────────────────────────────

\begin{definition}[Module-Topological Compression]
\label{def:compresion}
\textbf{Module-topological compression} is the process by which the
discrete Kuratowski phase entanglement forces non-planar monodromy
permutations onto the rotation system of a Causal Prism,
compelling the continuous Riemann surface (via Belyi's theorem)
to resolve the crossing by increasing its genus.

Formally: let $\mathcal{P}$ be a prism $K_{2,N}$ with
$g_{\mathrm{gen}} \ge 2$ phase classes. The phase anti-correlation
(Kuratowski pressure) entangles the cyclic edge order at the
poles. The entangled rotation system admits no planar face
decomposition. By Belyi's theorem
(Theorem~\ref{thm:belyi}), the unique Riemann surface that
resolves the dessin has genus $g = g_{\mathrm{gen}} - 1$.
\end{definition}

\noindent
\textbf{The mechanism chain.} The complete deductive sequence:
\begin{enumerate}
  \item The Poisson sprinkling creates the Hasse DAG (Axiom~1).
  \item The triangle-free topology forces prisms $K_{2,N}$.
  \item The phase trichotomy classifies the intermediaries into $\{+, 0, -\}$.
  \item At belly sizes $n \le n^*$, the Kuratowski constraints
        force phase anti-correlation.
  \item The anti-correlation entangles the monodromy permutation
        at the poles.
  \item The entangled monodromy fuses face pairs in
        $\sigma_\infty$.
  \item Fused faces increase the genus: $g =$ (number of entangled
        pairs).
  \item Belyi's theorem maps the dessin to a unique Riemann
        surface $X_g$.
  \item The Standard Model gauge fields are the Belyi functions
        $\beta \colon X_g \to \mathbb{P}^1$.
\end{enumerate}

\begin{result}[title={Mass IS geometry}]
Mass is not added to geometry---mass \textbf{is} geometry.
Each additional hole in the Riemann surface is a topological fold
forced by phase entanglement. The Standard Model mass hierarchy
is a genus staircase:
sphere → torus → double torus. The continuous gauge manifold is not
postulated---it is the obligatory resolution of discrete causal
constraints.
\end{result}

\stoneref{The Riemann surface is the \textbf{spectrum} of the
Heyting algebra of the Causal Prism. Its genus counts the number of
logically independent ``loops'' in the proof structure of the
prism---deductive paths that cannot be reduced to one another.
Genus~0 (sphere) means all proofs are homotopic: a single logical
channel ($U(1)$). Genus~1 (torus) means one independent loop: two
channels that cannot be deformed into each other ($SU(2)$ doublet).
Genus~2 (double torus) means two independent loops: three irreducible
channels ($SU(3)$ triplet). The gauge group is the \textbf{automorphism
group of the proof topology}.}

\begin{figure}[ht]
\centering
\begin{tikzpicture}[>=Stealth]
% Sphere (g=0)
\begin{scope}[shift={(-4,0)}]
  \draw[structure, thick, fill=structure!10] (0,0) circle (1.2);
  \draw[structure!40] (-1.2,0) arc (180:360:1.2 and 0.3);
  \draw[structure!40, dashed] (1.2,0) arc (0:180:1.2 and 0.3);
  \node[font=\small\bfseries, structure] at (0,1.8) {Sphere};
  \node[font=\small, second] at (0,-1.8) {$g=0$: $U(1)$};
  \node[font=\scriptsize, main] at (0,-2.3) {EM};
\end{scope}
% Torus (g=1)
\begin{scope}[shift={(0,0)}]
  \draw[third, thick] (0,0) ellipse (1.5 and 0.85);
  \draw[third, thick] (-0.5,0.1) arc (180:360:0.5 and 0.3);
  \draw[third, thick, dashed] (0.5,0.1) arc (0:180:0.5 and 0.3);
  \node[font=\small\bfseries, third] at (0,1.8) {Torus};
  \node[font=\small, second] at (0,-1.8) {$g=1$: $SU(2)$};
  \node[font=\scriptsize, main] at (0,-2.3) {Weak};
\end{scope}
% Double torus (g=2)
\begin{scope}[shift={(4.5,0)}]
  \draw[fourth, thick] (-0.8,0) ellipse (1.0 and 0.7);
  \draw[fourth, thick] (0.8,0) ellipse (1.0 and 0.7);
  \draw[fourth, thick] (-0.3,0.12) arc (180:360:0.3 and 0.2);
  \draw[fourth, thick, dashed] (0.3,0.12) arc (0:180:0.3 and 0.2);
  \draw[fourth, thick] (0.5,0.12) arc (180:360:0.3 and 0.2);
  \draw[fourth, thick, dashed] (1.1,0.12) arc (0:180:0.3 and 0.2);
  \node[font=\small\bfseries, fourth] at (0,1.8) {Double torus};
  \node[font=\small, second] at (0,-1.8) {$g=2$: $SU(3)$};
  \node[font=\scriptsize, main] at (0,-2.3) {Strong};
\end{scope}
\end{tikzpicture}
\caption{Gauge groups from genus. The Riemann surface of a
Generation~1 Prism is a sphere ($g=0$), whose monodromy gives
$U(1)$ (electromagnetism). Generation~2 lives on a torus ($g=1$),
producing $SU(2)$ (weak force). Generation~3 on a double torus
($g=2$) gives $SU(3)$ (strong force).}
\label{fig:gauge-genus}
\end{figure}

% ─────────────────────────────────────────────────────────────────────────────
\section{The Yang--Mills Continuum Limit}
\label{sec:yang-mills-continuum}
% ─────────────────────────────────────────────────────────────────────────────

The preceding sections established that the dessin's monodromy
associates a gauge group to each genus:
$g = 0 \to U(1)$, $g = 1 \to SU(2)$, $g = 2 \to SU(3)$.
The rigorous question remains: how is the \textbf{continuous Yang--Mills
Lagrangian} reconstructed from the permutation symmetry of the discrete
bipartite graph?

The answer uses the Grothendieck--Belyi correspondence already
introduced in Theorem~\ref{thm:belyi}, but now in the
reverse direction: from the dessin to the principal bundle.

\subsection{From $S_N$ to $SU(N)$: The Complete Derivation}

The promotion $S_N \to SU(N)$ is not an analogy: it is a chain
of four classical theorems, each rigorously established.
We present the complete derivation.

\medskip\noindent
\textbf{Starting data.}
A Causal Prism $K_{2,N}$ with $N$ intermediaries
$\{w_1, \ldots, w_N\}$ and permutations
$\sigma_0, \sigma_1 \in S_{2N}$ (the dessin structure). The
monodromy group is $G = \langle \sigma_0, \sigma_1 \rangle$, which
acts transitively on the $2N$~half-edges. The permutations
that reorder the intermediaries form a subgroup of $S_N$.

\medskip\noindent
\textbf{Step~I: From $S_N$ to the root system $A_{N-1}$
(Coxeter--Killing--Cartan).}

The symmetric group $S_N$ is a \textbf{Coxeter group} of
type~$A_{N-1}$: it is generated by the $N - 1$ simple transpositions
$s_i = (i,\; i{+}1)$ satisfying the relations
$s_i^2 = e$, $(s_i s_{i+1})^3 = e$, and $(s_i s_j)^2 = e$ for
$|i - j| > 1$. These relations uniquely determine the Dynkin
diagram~$A_{N-1}$~\cite{humphreys1990}:
\[
  \underset{s_1}{\circ} \text{---}
  \underset{s_2}{\circ} \text{---}
  \cdots \text{---}
  \underset{s_{N-1}}{\circ}
\]
The transpositions $s_i$ act on the space
$V = \{x \in \mathbb{R}^N : \sum x_j = 0\}$ by
$s_i(x) = x - (x \cdot \alpha_i)\,\alpha_i$ where
$\alpha_i = e_i - e_{i+1}$ are the simple roots. The complete
root system is $\Phi = \{e_i - e_j : i \neq j\}$.

\smallskip\noindent
\textbf{Uniqueness.} By the Killing--Cartan classification theorem,
the Dynkin diagram $A_{N-1}$ \textbf{uniquely} determines the
complex semisimple Lie algebra $\mathfrak{sl}(N, \mathbb{C})$, whose
compact real form is $\mathfrak{su}(N)$~\cite{humphreys1972}.

\smallskip\noindent
\textbf{Exclusion of alternatives.}
\begin{itemize}
  \item $O(N)$: its Weyl group is $(\mathbb{Z}/2)^{N-1} \rtimes S_N$
    (type $B_{N-1}$ or $D_{N-1}$), which contains reflections
    $x_i \mapsto -x_i$ absent from the pure permutation of
    intermediaries. $S_N$ alone does not generate $O(N)$.
  \item $U(N)$: $U(N) = SU(N) \times U(1)/\mathbb{Z}_N$. The
    central factor $U(1)$ would correspond to a global phase of
    the intermediaries. But $S_N$ (for $N \geq 3$) has trivial
    center: there is no permutation that commutes with all
    others and generates a continuous abelian factor.
  \item $Sp(N)$: its Weyl group is of type~$C_N$, with long
    and short roots; incompatible with the relations of $S_N$.
\end{itemize}

\medskip\noindent
\textbf{Step~II: From monodromy to flat connection
(Riemann--Hilbert).}

Belyi's Theorem (Theorem~\ref{thm:belyi}) produces a Riemann
surface~$X$ and a function $\beta \colon X \to \mathbb{P}^1$,
ramified over $\{0, 1, \infty\}$. The fundamental group
$\pi_1(\mathbb{P}^1 \setminus \{0, 1, \infty\})$ is free of
rank~2, generated by loops $\gamma_0, \gamma_1$ around $0$ and
$1$. The dessin's monodromy defines a representation:
\[
  \rho \colon \pi_1(\mathbb{P}^1 \setminus \{0, 1, \infty\})
  \;\longrightarrow\; S_N
  \;\hookrightarrow\; SU(N),
  \qquad
  \gamma_0 \mapsto \sigma_0, \quad \gamma_1 \mapsto \sigma_1,
\]
where the inclusion $S_N \hookrightarrow SU(N)$ sends each
permutation to the corresponding \textbf{permutation matrix}
in $U(N)$, which has determinant~$\pm 1$. For even permutations
($A_N \subset S_N$), the image falls in $SU(N)$; odd ones are
corrected by multiplying by $\det(\sigma)^{-1/N}$ ($N$-th root
of unity), absorbed into the global $U(1)$ phase which factors out.

By the \textbf{Riemann--Hilbert correspondence}~\cite{deligne1970},
the representation $\rho$ determines a unique (up to isomorphism)
flat vector bundle $(E, \nabla)$ of rank~$N$ over
$\mathbb{P}^1 \setminus \{0, 1, \infty\}$, with structure group
reduced to $SU(N)$ (the monodromy is unitary).

\medskip\noindent
\textbf{Step~III: From flat bundle to stable bundle
(Narasimhan--Seshadri).}

The Narasimhan--Seshadri theorem~\cite{narasimhan1965} establishes
a bijection between:
\begin{itemize}
  \item Irreducible unitary representations
    $\rho \colon \pi_1(X) \to SU(N)$.
  \item Stable holomorphic bundles of rank~$N$ and degree~0
    over~$X$.
\end{itemize}
The flat connection $\nabla$ from Step~II corresponds to a stable
holomorphic bundle $E \to X$ with structure group~$SU(N)$. This
is the \textbf{gauge principal bundle} from which the Yang--Mills
fields emerge.

\medskip\noindent
\textbf{Step~IV: From gauge bundle to Yang--Mills functional
(continuum limit).}

At finite density~$\rho$, the holonomy of $\nabla$ around an
elementary plaquette of the Hasse diagram (a causal quadrilateral of
area $\sim \rho^{-1/2}$) is an element of $SU(N)$ close to the
identity:
\[
  \mathrm{Hol}(\partial \square) = \mathcal{P}\exp\!\oint_{\partial \square}
  A \;\approx\; \mathbb{I} + F_{\mu\nu}\,\Delta x^\mu \Delta x^\nu
  + O(\rho^{-1}).
\]
The discrete action (sum over plaquettes of
$\mathrm{Tr}[\mathbb{I} - \mathrm{Hol}]$, as in Wilson's lattice
gauge theory~\cite{wilson1974}) converges, as
$\rho \to \infty$, to the continuous functional:

\begin{theorem}[Yang--Mills Continuum Limit]
\label{thm:yang-mills-continuum}
Let $K_{2,N}$ be a Causal Prism with intermediary permutation group
$S_N$, and let $X$ be the Riemann surface determined by the
associated Belyi dessin. In the continuum limit
$\rho \to \infty$:
\begin{enumerate}
  \item The root system $A_{N-1}$ of $S_N$ (Step~I) uniquely
    determines the Lie algebra $\mathfrak{su}(N)$.
  \item The Riemann--Hilbert correspondence (Step~II) produces a
    flat bundle with structure group $SU(N)$ over~$X$.
  \item The Narasimhan--Seshadri theorem (Step~III) identifies this
    bundle with a stable holomorphic bundle.
  \item Lattice convergence (Step~IV) produces the Yang--Mills
    functional:
\end{enumerate}
\begin{equation}
  S_{\mathrm{YM}} = \frac{1}{4g_{\mathrm{YM}}^2}
  \int_X \mathrm{Tr}(F_{\mu\nu} F^{\mu\nu})\,\sqrt{g}\,d^4x,
  \label{eq:yang-mills-action}
\end{equation}
where $F_{\mu\nu} = \partial_\mu A_\nu - \partial_\nu A_\mu
+ [A_\mu, A_\nu]$ is the curvature tensor of the $SU(N)$ principal
bundle.
\end{theorem}

\begin{proof}
Steps~I--IV constitute the proof.

(I)~$S_N$ is a Coxeter group of type $A_{N-1}$ (direct verification
of the Coxeter relations on the simple transpositions).
By the classification theorem for semisimple Lie algebras
(Killing~1888, Cartan~1894), the Dynkin diagram~$A_{N-1}$
uniquely determines~$\mathfrak{su}(N)$~\cite{humphreys1972}.

(II)~The dessin $(\sigma_0, \sigma_1)$ defines a representation
$\rho \colon \pi_1(\mathbb{P}^1 \setminus \{0,1,\infty\}) \to
S_N \hookrightarrow SU(N)$. Permutation matrices are
unitary, so the image of $\rho$ falls in $SU(N)$. By the
Riemann--Hilbert correspondence~\cite{deligne1970}, there exists a
unique flat vector bundle $(E, \nabla)$ of rank~$N$ with this
monodromy and structure group $SU(N)$.

(III)~Since $\rho$ is unitary and the monodromy group
$\langle \sigma_0, \sigma_1 \rangle$ acts transitively (by the
connected dessin hypothesis), the representation is irreducible.
By the Narasimhan--Seshadri theorem~\cite{narasimhan1965}, the
flat bundle is a stable holomorphic bundle over~$X$.

(IV)~The convergence of the Wilson lattice action to the continuous
Yang--Mills functional as the lattice spacing $a \to 0$ is a
standard result of lattice gauge
theory~\cite{wilson1974}: the plaquette action
$S_{\mathrm{plaq}} = \sum_{\square} \mathrm{Re}\,\mathrm{Tr}
[\mathbb{I} - U_\square]$ satisfies
$S_{\mathrm{plaq}} = \frac{a^4}{4g^2}\sum_\square
\mathrm{Tr}(F_{\mu\nu}F^{\mu\nu}) + O(a^6)$. In our framework,
$a \sim \rho^{-1/4}$ (Planck scale), so $\rho \to \infty$
is equivalent to $a \to 0$, and the discrete action converges to the
Yang--Mills functional~\eqref{eq:yang-mills-action}.
\end{proof}

\medskip\noindent
For the physically relevant case $N = 3$ (Generation~3), the
construction produces the gauge group $SU(3)$ (unique by
Killing--Cartan, Step~I). Where do exactly \textbf{eight}
gluons come from, if $|S_3| = 6$?

The three internal paths $w_1, w_2, w_3$ of $K_{2,3}$ form a
discrete Hilbert space of dimension~3: the fundamental representation
$\mathbf{3}$. A gauge boson is an operator that transitions a random
walker from one path to another; the space of all such transition
operators is the tensor product of a path and a dual path:
\[
  \mathbf{3} \otimes \bar{\mathbf{3}}
  \;=\; \mathbf{8} \oplus \mathbf{1}.
\]
The $\mathbf{1}$ component (trace, proportional to the identity) leaves
every path invariant and mediates no interaction: it is the color
singlet. The $3^2 - 1 = 8$ remaining traceless operators are the eight
Gell-Mann matrices---the eight gluon fields
$A_\mu^a$ ($a = 1, \ldots, 8$).

The permutation group $S_3$ with its $|S_3| = 6$ elements still
plays its role: it determines the \emph{type} of the Lie group (type
$A_2$, via Coxeter--Killing--Cartan), but the \emph{dimension} of
the boson space comes from $N^2 - 1$, not from~$|S_N|$.

\begin{itemize}
  \item \textbf{Confinement as topological inheritance}: the
    $K_5$ contraction (Theorem~\ref{thm:kuratowski-threat}) at the
    discrete level translates into the \textbf{confining flux tube}
    at the continuous level---the color field lines cannot spread
    because the planarity of the graph confines them.
\end{itemize}

\begin{result}[title={The 8 gluons from bipartite topology}]
\label{res:yang-mills-8-gluons}
The Yang--Mills Lagrangian with gauge group $SU(3)$ and its 8~gluon
fields is not postulated: it is \textbf{derived} from the $K_{2,3}$ prism.
The three internal paths form the fundamental representation
$\mathbf{3}$; the gauge bosons are the transition operators
$\mathbf{3} \otimes \bar{\mathbf{3}} = \mathbf{8} \oplus \mathbf{1}$;
the color singlet ($\mathbf{1}$) is discarded, leaving 8~gluons.
The group $S_3$ fixes the Lie type ($A_2 \to SU(3)$); the space of
transitions fixes the dimension ($3^2 - 1 = 8$). The gauge group is
$SU(N)$ and not $O(N)$, $U(N)$, or $Sp(N)$ because $S_N$ is a Coxeter
group of type~$A_{N-1}$---and no other.
\end{result}

\stoneref{In Stone duality, the promotion $S_N \to SU(N)$ is
the extension of a combinatorial symmetry (reordering of intermediate
lemmas in a proof) to a continuous symmetry (smooth rotation in proof
space). The three internal paths of $K_{2,3}$ are three intermediate
lemmas; the product $\mathbf{3} \otimes \bar{\mathbf{3}}$ is the space
of substitutions of one lemma for another. The trace ($\mathbf{1}$)
is the trivial substitution---leaving each lemma intact---and generates
no transformation. The 8~remaining directions are the 8~non-trivial
rotations of the internal structure of a proof with three lemmas.}
